\section*{Abstract}
\addcontentsline{toc}{section}{Abstract}
Declarative process models such as Dynamic Condition Response (DCR)
graphs specify the constraints an execution must satisfy rather than a
fixed control flow. This flexibility comes at a cost: every compliant
trace is treated as equally valid, with no mechanism to identify which
executions are operationally preferable in terms of cost, duration, or
resource usage. Existing solutions are either exact but unscalable
(constraint-based optimisers) or scalable but non-compliant (standard
reinforcement learning, which optimises reward without regard for the
process specification).

This thesis proposes a compliance-aware, multi-objective reinforcement
learning framework that closes this gap. A \ac{PPO} agent is trained
inside a DCR execution engine whose shield intercepts any non-enabled
action, leaving the marking unchanged; every trace generated during
training is therefore valid by construction, independently of reward
design or training budget. A normalised reward scalarises cost and
duration under configurable weights $(\alpha,\beta)$, so that pooling
the accepting traces of agents trained across a small weight grid
approximates the Pareto front of compliant executions. The framework is
further extended to a role-conditioned action space, where each event
can be executed by one of several resource roles, and to a
capacity-constrained variant that turns per-event role choice into a
genuine opportunity-cost allocation problem.

The framework is evaluated on two benchmark processes and externally
validated on independently mined DCR graphs never used during
development. It recovers the exact Pareto front on the smaller Expense
Report benchmark, and a well-populated, seed-stable 26-point front on
the larger Loan Application benchmark, where the role-conditioned
extension also learns event-level resource specialisation consistent
with the underlying cost--duration economics. External validation
replicates both findings on graphs of independent provenance, but also
exposes a genuine scalability ceiling: beyond roughly 42 activities
the unmasked learning agent's accept rate collapses, while the
shield alone, without learning, keeps producing valid executions well
past that point. These results show that compliance-by-construction
and multi-objective optimisation are not competing requirements, within
the scale at which the learning component remains effective.
